\newpage
\section{utils 模块}
\tkzname{utils} 模块提供底层的数值与解析工具，供 \tkzname{tkz-elements} 内部使用。这些函数保证了安全的数值转换，并为 TikZ 输出提供一致的字符串格式。


\subsection{\tkzname{utils} 模块函数一览表}

\begin{table}[htbp]
\centering
\caption{模块 \tkzMod{utils} 的函数。}

\begin{tabular}{@{}ll@{}}
\toprule
\texttt{函数} & \texttt{参考} \\
\midrule

\tkzFct{utils}{utils.parse\_point(str)}        & [\ref{sub:function_utils_parse_point}] \\

\tkzFct{utils}{utils.format\_number(r, n)}     & [\ref{sub:function_utils_format_number}]\\

\tkzFct{utils}{utils.format\_coord(x, decimals)} & [\ref{sub:function_utils_format_coord}]\\

\tkzFct{utils}{utils.format\_point(z, decimals)} & [\ref{sub:function_utils_format_point}]\\

\tkzFct{utils}{utils.checknumber(x, decimals)}  & [\ref{sub:function_utils_checknumber}]\\

\tkzFct{utils}{utils.almost\_equal(a, b, eps)}  & [\ref{sub:function_utils_almost_equal}]\\

\tkzFct{utils}{utils.clamp(x, min, max)} & [\ref{sub:function_utils_clamp}]\\

\tkzFct{utils}{utils.sign(x)}            & [\ref{sub:function_utils_sign}]\\

\tkzFct{utils}{utils.w(...)}             & [\ref{sub:function_utils_w}]\\

\tkzFct{utils}{utils.wlog(...)}              & [\ref{sub:function_utils_wlog}]\\
\bottomrule
\end{tabular}
\end{table}



\subsection{函数 \tkzFct{utils}{utils.parse\_point(str)}}
\label{sub:function_utils_parse_point}

解析形如 \code{"(x,y)"} 的字符串，并返回相应的数值坐标。该函数支持可选空格和科学计数法。

\medskip
\texttt{参数： } \code{str} – 表示一个点的字符串，例如 \code{"(3.5, -2.0)"}。

该函数接受：
\begin{itemize}
\item 十进制记法：\texttt{(3.5,-2)}
\item 科学计数法：\texttt{(1e-5,2E3)}
\item 可选空格
\item 可选括号
\end{itemize}

\medskip
\texttt{语法： }

\begin{center}
\code{local x, y = utils.parse\_point("(1.5, -2.3)")}
\end{center}

\medskip
\texttt{用途： }

该函数接收一个字符串参数并解析它，以数值形式取出 \code{x} 和 \code{y} 分量。输入字符串必须符合 \code{"(x, y)"} 的格式，其中 \code{x} 和 \code{y} 可以是用十进制或科学计数法书写的浮点数。


\medskip
\texttt{返回： }

两个数：\texttt{x}、\texttt{y} – 以 Lua 数表示的数值坐标。
两个数值：该点的实部和虚部。


\medskip
\texttt{用法示例： }

\begin{verbatim}
utils.parse_point("(3.5, -2)")
--  3.5   -2.0

utils.parse_point("(1e-5,2E3)")
-- 1e-05  2000

utils.parse_point(" -3.2e+1 , 4.5e-2 ")
-- -32  0.045
\end{verbatim}

科学计数法与十进制在数值上是等价的：

\begin{verbatim}
local x1 = utils.parse_point("(1e-5,0)")
local x2 = utils.parse_point("(0.00001,0)")
-- x1 == x2  --> true
\end{verbatim}

\medskip
\texttt{错误处理}

非法输入会引发 \texttt{tex.error}：

\begin{verbatim}
utils.parse_point("(a,b)")
-- Invalid point string
\end{verbatim}



\subsection{函数 \tkzFct{utils}{utils.format\_number(x, decimals)}}
\label{sub:function_utils_format_number}

该函数把一个数值（或数值字符串）格式化为固定小数位数的字符串表示。

\texttt{语法： }

\begin{center}
\code{local str = utils.format\_number(math.pi, 3)}
\end{center}

\texttt{用途： }

该函数把一个数（或可转换为数的字符串）转换成具有指定小数位数的字符串。在生成整洁的数值输出、或将其导出为 \TIKZ{} 坐标时，它尤其有用。

\texttt{参数： }

\begin{itemize}
\item \code{x} – 数值
\item \code{decimals} – 可选。小数位数（默认：5）
\end{itemize}

\texttt{返回： } 一个字符串，表示 \code{x} 具有指定小数位数的值。

\texttt{特点： }

\begin{itemize}
\item 尽可能自动把字符串转换为数。
\item 保证 \TIKZ{} 坐标或 LaTeX 输出的格式一致。
\item 输入不合法时引发错误。
\end{itemize}

\texttt{用法示例： }

\begin{verbatim}
utils.format_number(1e-5,5)
-- "0.00001"

utils.format_number(math.pi, 3)
-- "3.142"

utils.format_number("2.718281828", 2)
-- "2.72"

utils.format_number(2e3,2)
-- "2000.00"

utils.format_number(-4.2e-2,4)
-- "-0.0420"
\end{verbatim}


\subsection{在 tkz-elements 中的典型用法}

一个完整的解析与格式化流水线示例：

\begin{tkzexample}[latex=.4\textwidth]
\directlua{
init_elements()
local x,y = utils.parse_point("(1e-5,2e3)")

local sx = utils.format_number(x,5)
local sy = utils.format_number(y,2)
tex.print(sx)
tex.print('\\\\')
tex.print(sy)
}
\end{tkzexample}



\subsection{函数 \tkzFct{utils}{utils.format\_coord(x, decimals)}}
\label{sub:function_utils_format_coord}

该函数把一个数值格式化为固定小数位数的字符串。它是 \tkzFct{utils}{format\_number} 的一个轻量版本，用于输入已保证为数值的内部场合。

\texttt{语法： }

\begin{center}
\code{local s = utils.format\_coord(3.14159, 2)} \hfill $\rightarrow$ \code{"3.14"}
\end{center}

\texttt{参数： }

\begin{itemize}
\item \code{x} – 一个数（不做校验）。
\item \code{decimals} – 可选小数位数（默认：5）。
\end{itemize}

\medskip
\texttt{返回： }

一个具有固定小数格式的字符串。

\medskip
\texttt{说明： }

该函数由 \tkzFct{path}{add\_pair\_to\_path} 及其他路径构造方法在内部使用。

与 \tkzFct{utils}{format\_number} 不同，它不做输入校验，只应在已知输入为数值时使用。

\medskip
\texttt{相关函数： }

\begin{itemize}
\item \tkzFct{utils}{format\_number(x, decimals)} – 带校验的更安全的替代
\end{itemize}

\medskip
\texttt{精度说明： }

很小的数值可能根据小数位数被舍入为零：

\begin{verbatim}
utils.format_number(1e-8,5)
-- "0.00000"

utils.format_number(1e-8,10)
-- "0.0000000100"
\end{verbatim}

\texttt{重要提示： }

Lua 内部可能用科学计数法表示数：

\begin{verbatim}
print(1e-5)
-- 1e-05
\end{verbatim}

不过，\verb|utils.format_number| 总是生成适合 TikZ 输出的固定小数表示。

\subsection{与 \tkzname{path} 的交互}

在由计算所得的点构造 TikZ 路径时，格式化工具尤为重要。

\texttt{示例：路径中的科学计数法}

考虑一次计算后得到的一个非常小的坐标：

\begin{verbatim}
local x = 1e-5
local y = 2
\end{verbatim}

如果不做格式化就直接插入路径：

\begin{verbatim}
PA.curve:add_point(point(x,y))
\end{verbatim}

Lua 内部可能把该坐标表示为：

\begin{verbatim}
(1e-05,2)
\end{verbatim}

虽然数值正确，但科学计数法对 TikZ 输出并不理想。


\texttt{完整示例}

\begin{verbatim}
local p = path()
p:add_point(point(1e-5,0),5)
p:add_point(point(2e-5,1),5)

-- 生成：
-- (0.00001,0.00000) -- (0.00002,1.00000)
\end{verbatim}





\subsection{函数 \tkzFct{utils}{utils.checknumber(x, decimals)}}
\label{sub:function_utils_checknumber}

把一个数或数值字符串校验并转换为固定格式的小数字符串。

它是 \verb|utils.format_number| 的别名（为向后兼容而保留）。

\begin{verbatim}
utils.checknumber(1e-5,5)
-- "0.00001"
\end{verbatim}




\subsection{函数 \tkzFct{utils}{utils.format\_point(z, decimals)}}
\label{sub:function_utils_format_point}

把一个复点转换为适合坐标输出的字符串表示。

\texttt{语法： }

\begin{center}
\code{local s = utils.format\_point(z, 4)} \hfill $\rightarrow$ \code{"(1.0000,2.0000)"}
\end{center}

\medskip
\texttt{参数： }

\begin{itemize}
\item \code{z} – 一个带有 \code{re} 和 \code{im} 字段的表。
\item \code{decimals} – 可选精度（默认：5）。
\end{itemize}

\medskip
\texttt{返回： }

一个把该点表示为 \code{"(x,y)"} 的字符串。

\medskip
\texttt{错误处理。}

若 \code{z} 没有数值型的 \code{re} 和 \code{im} 分量，则引发错误。

\medskip
\texttt{相关函数： }

\begin{itemize}
\item \tkzFct{utils}{format\_coord}
\end{itemize}

\subsection{函数 \tkzFct{utils}{utils.almost\_equal(a, b, epsilon)}}
\label{sub:function_utils_almost_equal}

若两个数在给定容差内近似相等，则返回 \code{true}。

\medskip
\texttt{语法： }

\begin{center}
\code{if utils.almost\_equal(x, y) then ... end}
\end{center}

\medskip
\texttt{参数： }

\begin{itemize}
\item \code{a}、\code{b} – 待比较的两个数。
\item \code{epsilon} – 可选容差（默认：\code{tkz.epsilon}）。
\end{itemize}

\medskip
\texttt{返回： }

一个布尔值：若两值之差小于容差则为 \code{true}。



\subsection{函数 \tkzFct{utils}{utils.clamp(x, min, max)}}
\label{sub:function_utils_clamp}

返回把 \code{x} 限制在区间 $[\code{min}, \code{max}]$ 内的结果。

\medskip
\texttt{语法： }

\begin{center}
\code{local t = utils.clamp(t, 0, 1)}
\end{center}

\medskip
\texttt{返回： }

若 \code{x < min} 则返回 \code{min}；若 \code{x > max} 则返回 \code{max}；否则返回 \code{x}。

\subsection{函数 \tkzFct{utils}{utils.sign(x)}}
\label{sub:function_utils_sign}

返回一个数的符号。

\medskip
\texttt{语法： }

\begin{center}
\code{local s = utils.sign(z.A.im)}
\end{center}

\medskip
\texttt{返回： }

若 \code{x > 0} 则为 \code{1}；若 \code{x < 0} 则为 \code{-1}；若 \code{x = 0} 则为 \code{0}。

\subsection{函数 \tkzFct{utils}{utils.w(...)}}
\label{sub:function_utils_w}

把一条格式化消息写到终端和 .log 文件，并带有 \code{[tkz-elements]} 前缀。参数与 \code{string.format} 相同。若只想写到 .log 文件，请使用 \tkzFct{utils}{utils.wlog}。

\medskip
\texttt{语法： }

\begin{center}
\code{utils.w("Radius: \%.3f", C.OA.radius)}
\end{center}

\medskip
\texttt{返回： }

无返回值。

\subsection{函数 \tkzFct{utils}{utils.wlog(...)}}
\label{sub:function_utils_wlog}

只把一条格式化消息写入 .log 文件，并带有 \code{[tkz-elements]} 前缀。

\medskip
\texttt{语法： }

\begin{center}
\code{utils.wlog("Internal value: \%s", tostring(value))}
\end{center}

\medskip
\texttt{返回： }

无返回值。仅记录日志。
\endinput